\documentclass[10pt]{article}

\usepackage[margin=0.75in]{geometry}
\usepackage{amsmath,amsthm,amssymb,bm}
\usepackage{xcolor}
\usepackage{cancel}
\usepackage{graphicx}
\usepackage{changepage}
\usepackage{circuitikz}
\usepackage{pgfplots}
\usepackage{minted}
\usepackage{physics}
\usepackage{tikz}
\usepackage{tikz-3dplot}
\usepackage{hyperref}
\usepackage{siunitx}
\usepackage{fontspec}
\usepackage{relsize}
\usepackage{subfig}
\usepackage{todonotes}
\usepackage{contour}
\usepackage{multicol, multirow, booktabs}
\usepackage[breakable]{tcolorbox}
\usepackage[inline]{enumitem}

\theoremstyle{definition}
\newtheorem{problem}{Problem}
\newtheorem{soln}{Solution}

\pgfplotsset{compat=newest}
\usetikzlibrary{lindenmayersystems}
\usetikzlibrary{arrows}
\usetikzlibrary{calc}
\usetikzlibrary{positioning, fit}
\usetikzlibrary{3d, perspective}
\usetikzlibrary{math}
\usetikzlibrary{fadings}
\usetikzlibrary{angles,quotes} 

\definecolor{incolor}{HTML}{303F9F}
\definecolor{outcolor}{HTML}{D84315}
\definecolor{cellborder}{HTML}{CFCFCF}
\definecolor{cellbackground}{HTML}{F7F7F7}
\newcommand{\ui}{\hat{i}}
\newcommand{\uj}{\hat{j}}
\newcommand{\uk}{\hat{k}}
\newcommand{\ux}{\hat{\mathbf{x}}}
\newcommand{\uy}{\hat{\mathbf{y}}}
\newcommand{\uz}{\hat{\mathbf{z}}}
\newcommand{\uv}[1]{\hat{\bv{{#1}}}}
\newcommand{\pr}[1]{{#1^\prime}}
\newcommand{\pd}[2][]{{\frac{\partial^{#1}}{\partial {{#2}^{#1}}}}}
\newcommand{\dif}[2][]{{\frac{d^{#1}}{d{{#2}^{#1}}}}}
\newcommand{\grd}{\nabla}

\pgfdeclarelayer{background}  
\pgfsetlayers{background,main}
\AtBeginDocument{\RenewCommandCopy\qty\SI}

\makeatletter
\newcommand{\boxspacing}{\kern\kvtcb@left@rule\kern\kvtcb@boxsep}
\makeatother
\newcommand{\prompt}[4]{
    \ttfamily\llap{{\color{#2}[#3]:\hspace{3pt}#4}}\vspace{-\baselineskip}
}

\newcommand{\thevenin}[2]{
  \begin{center}
    \begin{circuitikz} \draw
      (0,0) -- (2,0) to[battery1, l_=$V_{Th}\eq#1$] (2,2) 
      to[resistor, l_=$R_{Th}\eq#2$] (0,2)
      ;
      \draw [o-] (-.07,2.079);
      \draw [o-] (-.07,0.079);
    \end{circuitikz}
  \end{center}
}

\newcommand{\norton}[2]{
  \begin{center}
    \begin{circuitikz} \draw
      (0,0) -- (3,0) to[american current source, l_=$I_{N}\eq#1$] (3,2) -- (0,2) (2,0)
      to[resistor, l=$R_{N}\eq#2$] (2,2)
      ;
      \draw [o-] (-.07,2.079);
      \draw [o-] (-.07,0.079);
    \end{circuitikz}
  \end{center}
}

\DeclareMathOperator{\Div}{div}
\DeclareMathOperator{\Curl}{curl}
\DeclareMathOperator{\sgn}{sgn}

\newcommand{\highlight}[1]{\colorbox{yellow}{$\displaystyle #1$}}
\newcommand{\ti}[1]{\widetilde{#1}}
\renewcommand{\div}[1]{\bv{\nabla}\cdot #1}
\renewcommand{\curl}[1]{\bv{\nabla}\times #1}

\newfontface{\Kaufmann}{Kaufmann}
\DeclareTextFontCommand{\kf}{\Kaufmann}
\newcommand{\scriptr}{\fontsize{12pt}{12pt}\kf{r}\fontsize{10pt}{10pt}}

\newfontface{\KaufmannB}{Kaufmann Bold}
\DeclareTextFontCommand{\kfb}{\KaufmannB}
\newcommand{\bscriptr}{\fontsize{12pt}{12pt}\kfb{r}\fontsize{10pt}{10pt}}
\newcommand{\justif}[2]{&{#1}&\text{#2}}
\newcommand{\bv}[1]{\bm{\mathrm{#1}}}

\title{Physics 3200Y: Assignment IX}
\author{Jeremy Favro (0805980) \\ Trent University, Peterborough, ON, Canada}
\date{\today}

\begin{document}
\maketitle

% PROBLEM 1
\begin{problem}
A metal bar of mass $m$ slides frictionlessly on two parallel conducting rails a distance $\ell$ apart.
A resistor $R$ is connected across the rails and a uniform magnetic field $\bv{B}$, pointing into the page,
fills the entire region.
\begin{enumerate}[label=(\alph*)]
  \item If the bar moves to the right with speed $v$, what is the current in the resistor? In what direction does it flow?
  \item What is the magnetic force on the bar? In what direction?
  \item If the bar starts out with speed $v_0$ at time $t=0$, and is left to slide, what is its speed at a later time $t$?
  \item The initial kinetic energy of the bar was, of course, $\frac{1}{2}mv_0^2$. Check that the energy delivered to the
        resistor is exactly $\frac{1}{2}mv_0^2$.
\end{enumerate}
\end{problem}
\begin{soln}~
  \begin{enumerate}[label=(\alph*)]
    \item We know that the current in a purely resistive circuit is related to the flux, and thereby the EMF, by
          $$-\frac{d\phi}{dt}=\epsilon=IR.$$
          Since the flux here is the field times the area of the loop, say the bar is at a position $x(t)$. Then
          $$\phi=B\ell x\implies IR=-B\ell\frac{dx}{dt}\implies I=-\frac{B\ell v}{R}.$$
          Where the negative sign indicates direction. Since the field is pointing into the page we can see that, using the
          RHR, the current is traveling downward through the resistor.
    \item Magnetic force is given by $\bv{f}=dq \bv{v}_{charges}\times \bv{B}$. Here since $\bv{v}_{charges}$ and $\bv{B}$ are at right
          angles to one another the cross product resolves to a regular product of magnitudes which we can integrate along the bar
          to find the force on the whole thing,
          $$\bv{F}=\int_{0}^{\ell}dq v_{charges} B(-\ux)=\int_{0}^{\ell}\frac{dq}{dt}Bd\ell'(-\ux)=I\ell B(-\ux)=-\frac{B^2\ell^2 v}{R}\ux.$$
          where we haven't picked up the sign on the current because we're considering magnitudes. This says that the force is leftward.
    \item The equation of motion here is
          $$m\ddot{x}=m\frac{dv}{dt}=-\frac{B^2\ell^2 v}{R}\implies v(t)=v_0\exp\left(-\frac{B^2\ell^2}{mR}t\right)$$
    \item The energy dissipated by the resistor is the power, given by $I^2R$:
          $$\frac{dW}{dt}=I^2 R=\frac{B^2\ell^2 v^2}{R}=\frac{B^2\ell^2}{R}v_0^2\exp\left(-2\frac{B^2\ell^2}{mR}t\right).$$
          Integrating this from $t=0$ to steady-state then
          \begin{align*}
            W & =\frac{B^2\ell^2}{R}v_0^2\int_0^\infty \exp\left(-2\frac{B^2\ell^2}{mR}t\right)dt                             \\
              & =-\frac{B^2\ell^2}{R}v_0^2\frac{mR}{2B^2\ell^2}\left[\exp\left(-2\frac{B^2\ell^2}{mR}t\right)\right]_0^\infty \\
              & =-\frac{1}{2}mv_0^2\left[e^{-\infty}-e^{0}\right]=\frac{1}{2}mv_0^2.
          \end{align*}
  \end{enumerate}
\end{soln}

% PROBLEM 2
\begin{problem}
An alternating current $I=I_0\cos\omega t$ flows down a long straight wire, and returns along a coaxial conducting tube
of radius $a$.
\begin{enumerate}[label=(\alph*)]
  \item In what \emph{direction} does the induced electric field point (radial, circumferential, longitudinal)?
  \item Assuming that the field goes to zero as $s \to \infty$, find $\bv{E}(s,t)$.
\end{enumerate}
\end{problem}
\begin{soln}
  \begin{enumerate}[label=(\alph*)]
    \item If we consider an Amp\`erian loop enclosing only the central conductor we see that the magnetic field
          is entirely in the $\uv{\phi}$ direction in cylindrical coordinates. Then since
          $$\curl{\bv{E}}=-\frac{d\bv{B}}{dt}$$
          $\curl{\bv{E}}$ must also be in $\uv{\phi}$. Looking at the $\uv{\phi}$ term for curl in
          cylindrical we see that
          $$\curl{\bv{v}}=\dots+\left[\frac{\partial v_s}{\partial z}-\frac{\partial v_z}{\partial s}\right]\uv{\phi}+\dots.$$
          Since the wire is long there shouldn't be any $z$ dependence so $\bv{E}$ should only have a $\uz$ component which varies with
          $s$.
          \todo{Contour diagrams\dots}
    \item Taking a new contour of length $\ell$ which slices the outer conductor we note that using our definition of flux,
          $$\int (\curl{\bv{E}})\cdot d\bv{a}=\oint \bv{E}\cdot d\ell'\uz=-\frac{d}{dt}\int \bv{B}\cdot d\bv{a}$$
          where the first equality comes from Stoke's (curl) theorem. We can also use Ampere's law to write
          $$\oint \bv{B}\cdot d\bv{\ell}'=2\pi s' B=\mu_0 I_0\cos\omega t.$$
          Plugging this into the chain of equalities for the electric field gives that
          $$E\ell=-\frac{d}{dt}\int_0^\ell\int_s^a \frac{\mu_0 I_0}{2\pi s'}\cos\omega t ds'dz'$$
          and so
          $$E=\frac{\mu_0 I_0 \omega}{2\pi}\ln\left(\frac{a}{s}\right)\sin\omega t.$$
  \end{enumerate}
\end{soln}

% PROBLEM 3
\begin{problem}
An infinite cylinder of radius $R$ carries a uniform surface charge $\sigma$.
We propose to set it spinning about its axis, at a final angular velocity $\omega_f$. How much work will this take per unit length?
Do it two ways, and compare your answers:
\begin{enumerate}[label=(\alph*)]
  \item Find the magnetic field and the induced electric field (in the quasistatic approximation),
        inside and outside the cylinder, in terms of $\omega$, $\dot{\omega}$, and $s$. Calculate the torque you must
        exert, and from that obtain the work done per unit length ($W=\int Nd\phi$).
  \item Use Eq. 7.44 ($W=\frac{1}{2\mu_0}\int_{\mathbb{R}}B^2d\tau$) to determine the energy stored in the resulting magnetic field.
\end{enumerate}
\end{problem}
\begin{soln}~
  \begin{enumerate}[label=(\alph*)]
    \item When we set this system rotating it becomes a solenoid. We know the field of a solenoid to be
          $$\bv{B}=\begin{cases}
              \mu_0\sigma\omega R\uz & s<R  \\
              0                      & s> R
            \end{cases}$$
          where $\sigma\omega R$ is the surface current density. Since, for a contour with area in $\uz$
          $$\oint \bv{E}\cdot d\bv{\ell}'=-\frac{d}{dt}\int\bv{B}\cdot d\bv{a}$$
          and we expect that the electric field is entirely in $\uv{\phi}$, as is $d\bv{\ell}'$,
          $$E2\pi s=-\frac{d}{dt}\int_{0}^{s}s'ds'\int_{0}^{2\pi}d\phi' \mu_0\sigma\omega R=-2\pi\mu_0\sigma\dot{\omega}\frac{R s^2}{2}$$
          so
          $$\bv{E}=-\mu_0\sigma\dot{\omega}\frac{R s}{2}\uv{\phi}$$
          inside. To do outside we change the limits of the integrals,
          $$E2\pi s=-\frac{d}{dt}\int_{0}^{R}s'ds'\int_{0}^{2\pi}d\phi' \mu_0\sigma\omega R=-2\pi\mu_0\sigma\dot{\omega}\frac{R^3}{2}$$
          so
          $$\bv{E}=-\mu_0\sigma\dot{\omega}\frac{R^3}{2s}$$
          outside. At the surface then (I think) the electric field is
          $$\bv{E}=-\mu_0\sigma\dot{\omega}\frac{R^2}{2}\uv{\phi}.$$
          Torque then is
          $$\bv{\tau}=\bv{r}\times\bv{F}=R(2\pi R\ell \sigma)(-\mu_0\sigma\dot{\omega}\frac{R^2}{2})\uz$$
          since the total charge on a segment of length $\ell$ is $2\pi R\ell \sigma$.
          This works out to
          $$\bv{\tau}=-\pi\ell\mu_0\sigma^2\dot{\omega}R^4\uz.$$
          Work per $\ell$ then is 
          $$\frac{W}{\ell}=-\pi\mu_0\sigma^2R^4\int\dot{\omega}d\phi$$
          but 

    \item Use Eq. 7.44 ($W=\frac{1}{2\mu_0}\int_{\mathbb{R}}B^2d\tau$) to determine the energy stored in the resulting magnetic field.
  \end{enumerate}
\end{soln}
\end{document}